\documentclass[12pt,a4paper]{article}

\usepackage[T2A]{fontenc}
\usepackage[utf8]{inputenc}
\usepackage[russian]{babel}
\usepackage{amsmath,amssymb,amsfonts}
\usepackage{geometry}
\usepackage{booktabs}
\usepackage{enumitem}
\usepackage{array}
\usepackage{hyperref}

\geometry{margin=2.2cm}

\title{IISVersion7 и оператор Sumigron}
\author{Проект анализа интегрального индекса состояния}
\date{\today}

\newcommand{\sgsumintright}[1]{\mathop{\mathchoice
  {\ooalign{
    \hfil$\displaystyle\sum$\hfil\cr
    \hfil$\displaystyle\int$\hfil\cr
    \hidewidth\raisebox{-1.35ex}{\kern1.25ex$\scriptstyle #1$}\hidewidth\cr
  }}
  {\ooalign{
    \hfil$\textstyle\sum$\hfil\cr
    \hfil$\textstyle\int$\hfil\cr
    \hidewidth\raisebox{-1.00ex}{\kern1.00ex$\scriptscriptstyle #1$}\hidewidth\cr
  }}
  {\ooalign{
    \hfil$\scriptstyle\sum$\hfil\cr
    \hfil$\scriptstyle\int$\hfil\cr
    \hidewidth\raisebox{-0.75ex}{\kern0.75ex$\scriptscriptstyle #1$}\hidewidth\cr
  }}
  {\ooalign{
    \hfil$\scriptscriptstyle\sum$\hfil\cr
    \hfil$\scriptscriptstyle\int$\hfil\cr
    \hidewidth\raisebox{-0.60ex}{\kern0.55ex$\scriptscriptstyle #1$}\hidewidth\cr
  }}
}}

\newcommand{\sgop}{\sgsumintright{\bowtie}}
\newcommand{\sigmoid}{\sigma}

\begin{document}
\maketitle

\section{Назначение версии V7}
IISVersion7 вводится как следующая теоретическая версия после V5 и V6.
Её главная идея состоит в том, чтобы:
\begin{itemize}[leftmargin=1.25cm]
    \item не сжимать окно сигнала слишком рано одной ограниченной нелинейностью;
    \item заменить плоскую раннюю агрегацию окна на структурный оператор \textbf{Sumigron};
    \item сохранить нелинейные компоненты \(A,\Gamma,V,Q\) и gated-логику итогового индекса;
    \item сохранить трёхмерную интерпретацию состояния через \((IIS, RES, DYN)\).
\end{itemize}

Иными словами, V7 строится как иерархия:
\[
\text{окно} \;\to\; \sgop \;\to\; A,\Gamma,V,Q \;\to\; IIS,RES,DYN.
\]

\section{Оператор Sumigron}
\subsection{Канонический знак}
Для оператора Sumigron фиксируется знак
\[
\sgop(X,M),
\]
где:
\begin{itemize}[leftmargin=1.25cm]
    \item \(\sum\) обозначает агрегацию по структуре окна;
    \item \(\int\) обозначает непрерывный проход по форме сигнала;
    \item маленький индекс \(\bowtie\) справа снизу обозначает структурное слияние нескольких характеристик окна.
\end{itemize}

В строгой резервной нотации можно писать также
\[
\mathcal{S}_{\mathrm{sg}}(X,M).
\]

\subsection{Вход оператора}
Пусть \(W_{i,c}=\{x_{i,c,t}\}_{t=1}^{L_i}\) --- окно длины \(L_i\) для окна \(i\) и канала \(c\).

Сначала считаются обычные статистики:
\[
\mu_{i,c}=\frac{1}{L_i}\sum_{t=1}^{L_i}x_{i,c,t},
\]
\[
\sigma_{i,c}=\sqrt{\frac{1}{L_i}\sum_{t=1}^{L_i}(x_{i,c,t}-\mu_{i,c})^2},
\]
\[
\hat{x}_{i,c,t}=\frac{x_{i,c,t}-\mu_{i,c}}{\sigma_{i,c}+\varepsilon}.
\]

\subsection{Attentive-ядро}
Внутри окна вводится локальный драйв внимания:
\[
u_{i,c,t}=0.60\,\hat{x}_{i,c,t}+0.40\,|\hat{x}_{i,c,t}|.
\]

Далее вычисляются attention-веса:
\[
\alpha_{i,c,t}=
\frac{\exp(u_{i,c,t}/\tau)}
{\sum_{s=1}^{L_i}\exp(u_{i,c,s}/\tau)}.
\]

На их основе задаются три структурные компоненты окна:
\[
\mu^{sg}_{i,c}=\sum_{t=1}^{L_i}\alpha_{i,c,t}\,x_{i,c,t},
\]
\[
\sigma^{sg}_{i,c}=
\sqrt{
\sum_{t=1}^{L_i}\alpha_{i,c,t}
\bigl(x_{i,c,t}-\mu^{sg}_{i,c}\bigr)^2
},
\]
\[
E^{sg}_{i,c}=
\log\!\left(
1+\sum_{t=1}^{L_i}\alpha_{i,c,t}\,x_{i,c,t}^2
\right).
\]

\subsection{Скалярная форма Sumigron}
После центрирования и нормировки этих трёх каналов:
\[
\tilde{\mu}^{sg}_{i,c},\qquad
\tilde{\sigma}^{sg}_{i,c},\qquad
\tilde{E}^{sg}_{i,c},
\]
Sumigron определяется как
\[
\sgop(W_{i,c})=
w_\mu \tilde{\mu}^{sg}_{i,c}
+w_\sigma \tilde{\sigma}^{sg}_{i,c}
+w_E \tilde{E}^{sg}_{i,c}.
\]

В базовой теоретической настройке можно брать:
\[
w_\mu=0.32,\qquad
w_\sigma=0.28,\qquad
w_E=0.40.
\]

\section{Компоненты IISVersion7}
\subsection{Эмоционально-корковый компонент \(A_7\)}
\[
A_7=
\tanh\!\left(
a_1\,\sgop(W^{\alpha asym}_i)
+a_2\,\sgop(W^{total\ asym}_i)
\right).
\]

\subsection{Gamma-компонент \(\Gamma_7\)}
\[
\Gamma_7=
\tanh\!\left(
g_1\,\sgop(W^\gamma_i)
+g_2\,\sgop(W^{\gamma/\alpha}_i)
\right).
\]

\subsection{Вегетативный компонент \(V_7\)}
\[
V_7=
\tanh\!\left(
v_1\,\sgop(W^{HR}_i)
+v_2\,\sgop(W^{HF/LF}_i)
\right).
\]

\subsection{Интегральное качество \(Q_7\)}
Линейное ядро:
\[
Q^{lin}_7=q_1A_7+q_2V_7-q_3\Gamma_7.
\]

Когерентность:
\[
Q^{coh}_7=A_7V_7-\lambda_1|A_7-V_7|-\lambda_2\max(\Gamma_7,0).
\]

Итоговая формула:
\[
Q_7=
\tanh\!\left(
\eta_1Q^{lin}_7+\eta_2Q^{coh}_7
\right).
\]

\section{Итоговый индекс IISVersion7}
\subsection{Gated-логиты режимов}
\[
z^{reg}_7=b_1Q_7+b_2V_7-b_3\Gamma_7+b_4A_7,
\]
\[
z^{mob}_7=m_1V_7-m_2\Gamma_7-m_3Q_7+m_4A_7,
\]
\[
z^{dep}_7=d_1\Gamma_7-d_2V_7-d_3Q_7-d_4A_7.
\]

Мягкие веса:
\[
\bigl(g^{reg}_7,g^{mob}_7,g^{dep}_7\bigr)
=
\operatorname{softmax}(z^{reg}_7,z^{mob}_7,z^{dep}_7).
\]

\subsection{Локальные режимные оценки}
\[
s^{reg}_7=\sigmoid(c_1Q_7+c_2V_7-c_3\Gamma_7),
\]
\[
s^{mob}_7=\sigmoid(r_1A_7+r_2V_7-r_3\Gamma_7),
\]
\[
s^{dep}_7=\sigmoid(-p_1V_7-p_2Q_7-p_3A_7+p_4\Gamma_7).
\]

\subsection{Финальная формула}
\[
IIS_7=
g^{reg}_7 s^{reg}_7
+g^{mob}_7 s^{mob}_7
+g^{dep}_7 s^{dep}_7.
\]

\section{Дополнительные координаты трёхмерной карты}
\subsection{Ресурсность}
\[
RES_7=
\operatorname{clip}\!\left(
0.5+0.5\tanh\!\bigl(
\rho_1V_7-\rho_2\max(\Gamma_7,0)-\rho_3|A_7-V_7|
\bigr),
0,1
\right).
\]

\subsection{Динамика}
\[
DYN_7=
\tanh\!\left(
\delta_1(IIS_{7,i}-IIS_{7,i-1})
+\delta_2(RES_{7,i}-RES_{7,i-1})
-\delta_3\,Vol_i
\right).
\]

Таким образом, трёхмерная точка состояния имеет вид
\[
Z_i=(IIS_{7,i},RES_{7,i},DYN_{7,i}).
\]

\section{Теоретические преимущества версии V7}
\subsection{Преимущества относительно раннего сжатия через \texorpdfstring{\( \tanh \)}{tanh}}
\begin{enumerate}[leftmargin=1.25cm]
    \item \textbf{Сохранение структуры окна.} В V7 окно не схлопывается сразу в один усреднённый bounded-ответ. Сначала сохраняются уровень, вариативность и энергия.
    \item \textbf{Сохранение амплитуды.} Нелинейность переносится выше по иерархии. Поэтому амплитуда и внутренняя геометрия окна теряются меньше.
    \item \textbf{Лучшая работа со структурными режимами.} Если состояния отличаются не только средним уровнем, но и распределением пиков, шумом или burst-поведением, Sumigron теоретически должен различать их лучше.
\end{enumerate}

\subsection{Преимущества самого оператора Sumigron}
\begin{enumerate}[leftmargin=1.25cm]
    \item \textbf{Разделение каналов информации.} Sumigron хранит отдельно \(\mu\), \(\sigma\) и \(E\), а не прячет их в одном числе.
    \item \textbf{Attention без потери интерпретации.} Внимание применяется к локальной структуре окна, но итог остаётся разложимым на понятные каналы.
    \item \textbf{Масштабируемость.} Один и тот же оператор можно применять к EEG, ECG, HRV, EDA и другим модальностям.
    \item \textbf{Совместимость с маркировкой.} Sumigron естественно допускает typed-расширение \(\sgop(X,M)\), где \(M\) кодирует тип сигнала, происхождение и контекст.
\end{enumerate}

\subsection{Преимущества IISVersion7 как системы}
\begin{enumerate}[leftmargin=1.25cm]
    \item \textbf{Иерархичность.} V7 естественно собирает в одну цепочку окно, оператор, компоненты, gated-индекс и 3D-интерпретацию.
    \item \textbf{Гибкость.} Можно заменить только ранний агрегатор окна, не ломая компоненты \(A,\Gamma,V,Q\) и не уничтожая совместимость с V5/V6.
    \item \textbf{Более богатая геометрия состояния.} При сохранении \((IIS,RES,DYN)\) V7 остаётся совместимой с 2D и 3D картой состояний.
    \item \textbf{Лучший теоретический потенциал для 3+ режимов.} За счёт менее плоской ранней агрегации V7 потенциально лучше подходит для выделения промежуточных и переходных состояний.
\end{enumerate}

\section{Ограничения}
IISVersion7 в данной записи является \textbf{теоретической} архитектурой.
Это означает:
\begin{itemize}[leftmargin=1.25cm]
    \item формулы задают осмысленную структуру, но ещё не доказывают выигрыш на реальных данных;
    \item коэффициенты \(a_1,a_2,\dots,\delta_3\) должны подбираться и проверяться на реальных датасетах;
    \item выигрыш V7 должен подтверждаться отдельно по state-capacity, stability и dynamic-метрикам.
\end{itemize}

Самая точная интерпретация на текущем этапе:
\[
\text{V7 = теоретически более богатая иерархическая версия, но ещё не эмпирически финальная.}
\]

\end{document}
